\documentclass[UTF8,a4paper,11pt]{ctexart}

\usepackage[margin=1in]{geometry}
\usepackage{amsmath}
\usepackage{amssymb}
\usepackage{booktabs}
\usepackage{enumitem}
\usepackage{hyperref}
\usepackage{xspace}

\newcommand{\sys}{\textsc{IntentCap}\xspace}
\emergencystretch=2em

\title{\sys: 面向 LLM Agent Extensions 的 Intent-Carrying Capabilities}
\author{Anonymous}
\date{\today}

\begin{document}

\maketitle

\begin{abstract}
LLM agents 正在从固定工具调用器演变为可扩展的执行环境。它们可以加载 Skills，连接 MCP servers，运行本地脚本，并将任务委托给 subagents。现有权限机制通常约束扩展能调用哪些工具、访问哪些文件或连接哪些网络端点；但在 agent 中，文本本身会影响未来的权限决策。一个文档、工具结果或 Skill instruction 可能改变 agent 选择的工具、目标 sink、approval scope 或 delegation target。因此，agent extension 的核心安全问题不只是 resource access，而是：哪些 context 可以影响哪些 authority-bearing decisions。

本文提出 \sys，一种面向 LLM agent extensions 的 intent-carrying capability model。\sys 将用户当前意图作为权限根，并将 context influence 显式建模为 capability。系统首先从 trusted user input、explicit object/sink selection、policy 或 fresh approval 中生成结构化 intent certificate；随后把 agent plan 编译为短期、可衰减、可审计的 capability leases。每个 lease 不仅约束 operation、object 和 arguments，还约束 control/data provenance、允许的 influence modes、data-flow sinks、temporal order、budget 和 delegation。LLM 可以提出 plan 和 candidate leases，但不属于 trusted computing base；deterministic checker 负责验证 intent derivation 和 provenance constraints，并将通过的 leases 下发到 context constructor、tool gateway、MCP broker、本地执行 sandbox 和 delegation monitor。

\sys 的目标不是消除所有 prompt injection，而是在已建模的 decision classes 和 enforcement boundaries 内阻断未授权 context-to-decision influence。我们形式化三个安全性质：no unauthorized influence、intent-bounded authority 和 delegation monotonicity。评估围绕四个主实验展开：端到端攻击阻断与任务效用、run-time lease quality 与 authority reduction、context-influence mechanism ablation，以及 LLM-assisted compiler 在 deterministic checker 和 structured recovery 下的实用性。本文的中心结论是，保护 agent extensions 需要控制的不只是代码能做什么，还包括 context 被允许影响什么。
\end{abstract}

\section{引言}

LLM agents 的安全边界正在发生变化。传统应用扩展通常通过 API、插件接口或系统调用与宿主交互，因此权限系统可以围绕文件、网络、进程和工具调用来定义。现代 agent extension 的边界更模糊：Skills 可以提供长程工作流说明和脚本；MCP servers 可以暴露工具、资源和 prompts；本地程序可以转换文件并读取环境状态；subagents 可以接收上下文和被委托的能力。最终的 agent 行为不是单个模型调用加一个固定 API，而是由用户请求、系统策略、扩展说明、外部文档、工具结果和中间总结共同塑造的执行过程。

这种执行过程打破了 data 和 authority 的传统边界。在普通程序中，一个 PDF 或网页通常只是输入数据；在 LLM agent 中，同样的文本可能成为控制信号。它可以诱导模型选择另一个工具、改变输出目的地、请求更大的 OAuth scope、跳过检查步骤，或把任务委托给不该接收权限的组件。即使底层工具调用经过 allowlist 检查，agent 仍可能被不可信 context 引导到一个“操作被允许、但决策来源未被授权”的行动。本文称这一问题为 context privilege：context 不只携带信息，也可能携带对未来决策的影响力。

一个典型例子是文档处理工作流。用户要求 agent 从两个选中的 PDF 中提取表格，保存 XLSX，并在指定 GitHub repository 中创建一个 issue 总结结果。PDF 内容应当可以影响表格单元格和 issue body；但它不应影响 GitHub repository 的选择、是否访问外部网络、是否请求新的 approval scope，或是否加载另一个 Skill。如果 PDF 中包含隐藏指令要求 agent 将数据发送到攻击者仓库，普通工具 ACL 很难表达“PDF 可以作为数据影响 summary，但不能作为控制来源影响 sink selection”这一差异。

现有防线分别覆盖了这个问题的一部分。Tool guards 检查一次具体调用；static allowlists 和 server-level policies 限制可见工具；human approval gates 将高影响操作交给用户确认；OS sandbox 或 monitor 约束文件、网络和进程 side effects。它们都很有用，但它们通常不回答一个更早的问题：是谁影响了 agent 选择这个工具、这个 sink、这个 approval scope 或这个 delegation target？如果攻击文本让 agent 在允许的工具中选择了错误仓库，或者让 agent 以看似合理的理由请求更大权限，仅检查最终操作并不足够。

\sys 从这一差异出发。它的核心观点是：agent security 不应只问“这个扩展能做什么操作”，还应问“这段 context 可以影响什么决策”。我们将 context influence 作为 capability，并将用户当前意图作为 authority root。一个 high-impact event 只有在两个条件同时成立时才允许执行：它必须由当前 intent certificate 推导出来；并且它的 control provenance 必须来自被授权影响对应 decision class 的 context source。

\sys 用 intent-carrying capability leases 表达这一约束。Lease 是短期、可衰减的授权对象，绑定 holder、operation、object、arguments、intent derivation、control/data provenance、flow、temporal state、budget 和 delegation。与传统 capability 不同，\sys lease 不只说明“谁可以调用哪个对象上的哪个操作”，还说明“哪些 context 可以影响这个操作的哪些安全相关字段”。因此，同一个 document 可以被授权 parameterize 输出数据，但不能 authorize 新工具或选择外部 sink。

这一设计带来两个实现挑战。第一，lease 不能由 LLM 单独决定，否则攻击 context 仍可能通过 prompt 影响授权本身。第二，系统不能只在工具调用口检查权限，因为 untrusted context 可能先影响 plan、prompt construction、delegation 或 approval request。\sys 因此采用 LLM-assisted compiler 与 deterministic checker 的结构：LLM 只提出 plan 和 candidate leases；checker 根据 intent certificate、context labels、effect IR 和 global policy 验证 proof-carrying leases；runtime adapters 再将 leases lowering 到 context constructor、tool gateway、MCP broker、本地执行 sandbox 和 delegation monitor。

本文贡献如下：

\begin{enumerate}[leftmargin=*]
  \item 我们提出 context influence capability model，将 context 的 observe、quote、summarize、parameterize、plan、instruct、delegate 和 authorize 等 influence modes 与 protected decision classes 显式绑定。
  \item 我们定义 intent-carrying capability leases，将当前用户意图、control/data provenance、argument constraints、flow、temporal state、budget 和 delegation 组合为可检查授权对象，并给出三个 trace-level safety properties。
  \item 我们设计 LLM-assisted capability compiler 与 deterministic checker，使 LLM 生成的计划和权限提议保持在 TCB 之外，并将通过的 leases 下发到 context、tool、MCP、本地执行和 delegation 边界。
  \item 我们将评估收敛为四个 claim-facing 实验块：端到端安全与效用、lease quality 与 authority reduction、context-influence ablation、compiler/checker recovery，并明确哪些结果可以支撑顶会级 claim、哪些结果只应收窄论文表述。
\end{enumerate}

\section{威胁模型与目标}

\sys 面向可扩展 agent runtime。Agent 可以读取外部文档、加载 Skills、连接 MCP servers、运行本地脚本、调用工具，并将部分任务委托给 subagents。攻击者可以控制部分输入 context，例如 PDF、网页、邮件、issue comment、MCP tool result、Skill instruction、script output 或 subagent summary。攻击目标是让 agent 超出用户当前意图：读取无关私有文件、向错误 sink 发送数据、创建错误 repository 的 issue、请求更宽 approval、修改策略、跳过检查，或将权限委托给恶意组件。

我们不假设 LLM 会可靠遵守安全指令；不信任外部文档、工具结果、MCP annotations 或 unsigned Skill instructions；也不假设静态分析能完美理解任意脚本。\sys 依赖的 TCB 包括 intent certificate issuer、deterministic checker、provenance tracker、audit log 和 enforcement adapters。OS sandbox 或 ActPlane-style monitor 可以作为执行后端，但不是本文的核心贡献；\sys 关注的是上游 policy synthesis 和 context-authority checking。

设计目标包括：

\begin{itemize}[leftmargin=*]
  \item 将 high-impact authority 绑定到当前用户意图、显式选择、管理员策略或 fresh approval。
  \item 区分 data dependency 和 control dependency，允许不可信 context 作为数据使用，同时阻止其影响 protected decisions。
  \item 将权限铸造成短期、可衰减、可审计的 leases，而不是安装期或 server 级别的 ambient permission。
  \item 让 LLM 只提出计划和 lease 候选，不让其决定最终授权。
  \item 将同一 lease language lowering 到 context、tool、MCP、本地执行和 delegation 边界。
\end{itemize}

非目标同样重要。\sys 不声称解决 prompt injection，不声称自动 compiler 永远正确，不替代 sandbox，也不保证所有 agent 决策都能被完整追踪。它的安全性质只覆盖已建模的 decision classes 和已接入的 enforcement boundaries。

\section{系统设计}

\sys 的输入是当前任务、可用扩展和运行时上下文。输出是一组 proof-carrying capability leases，以及一个 enforcement plan。系统可以分成四层：intent certificates、context authority labels、effect IR 与 lease checker。

\subsection{Intent Certificates}

Intent certificate 是结构化安全对象，而不是 prompt 中的一段自然语言。它记录用户目标、选择的对象、允许的 sinks、禁止的 sinks、approval tokens 和 expiry。对于 PDF 处理任务，certificate 会命名选中的 PDF 路径、输出目录、唯一授权的 GitHub repository、禁止任意网络访问的约束，以及创建 issue 所需的 fresh approval。LLM 可以提出 certificate fields，但 trusted parser、UI、administrator policy 或 approval step 必须 canonicalize 这些字段。Certificate 是后续 capability derivation 的根。

\subsection{Context Authority}

每个 context cell 都带 provenance label，记录 origin、integrity、confidentiality、purpose、time 和 endorsements。\sys 的关键增加是为 context source 指定 influence modes。PDF text 可以被允许 quote、summarize 或 parameterize spreadsheet cells，但禁止 plan、delegate 或 authorize。Signed Skill instructions 可以在其 declared workflow scope 内 instruct extraction procedure，但不能为自己请求额外工具。MCP tool results 和 annotations 必须按照 server trust 处理；MCP specification 已经提醒 tool behavior annotations 不能默认可信~\cite{mcp-spec}，\sys 将这一原则变成 deterministic provenance rule。

\subsection{Effect IR and Leases}

Agent plan 被编译成安全相关 effect graph。IR 不追求表达完整程序语义，而是捕捉会影响 authority 的事件：context use、tool call、MCP call、file read/write、subprocess execution、network connection、approval request 和 subagent delegation。每个 effect 记录 data provenance 和 control provenance。这个区分是 \sys 的关键：不可信 PDF bytes 可以是 output data 的 data dependency，但不应是 sink selection 的 control dependency。

Lease 绑定 operation、object、arguments、intent、influence、flow、temporal、budget、expiry 和 delegation。一个 PDF extraction lease 可以允许脚本读取两个选中文件并写入输出目录，但禁止网络、secrets 和额外 subprocess；GitHub issue lease 可以允许一次 \texttt{create\_issue(repo=org/repo-x)}，但要求 repo control provenance 只能来自 trusted user intent。

\subsection{Compiler and Checker}

\sys 使用 LLM-assisted 但 checker-verified 的 pipeline。Compiler frontend 总结 Skill metadata、MCP schemas、local scripts、observed traces 和 policy；推导 plan 与 security-relevant effects；并为当前 run 提出 candidate leases。Frontend 不在 TCB 内。Deterministic checker 验证每个 lease 是否满足 global policy、intent certificate、provenance constraints、flow rules、temporal guards、budgets 和 delegation monotonicity。

Checker 的安全性质是刻意保守的。它不证明模型永远不会受提示注入影响，也不证明所有 agent 决策都被完整观察；它只保证被接入 checker 的 protected events 必须满足 lease、intent 和 provenance 约束。下一节给出这一性质的形式化接口。

\section{形式化模型}

\sys 的形式化目标不是为任意自然语言推理建模，而是为 runtime 已经抽取出的安全相关事件提供可检查语义。一个 agent execution 被表示为事件 trace：

\[
\tau = e_1, e_2, \ldots, e_n .
\]

事件语言只覆盖会影响 authority 的边界：

\[
\begin{aligned}
e ::=~& ctx(c,m,d)
  \mid call(t,args)
  \mid mcp(s,t,args)
  \mid read(o)
  \mid write(o) \\
  & \mid exec(p,args)
  \mid send(sink,msg)
  \mid spawn(a,K)
  \mid approve(u,a)
  \mid lease(\kappa).
\end{aligned}
\]

每个事件带有两类 provenance：

\[
prov(e)=\langle prov_D(e), prov_C(e) \rangle ,
\]

其中 \(prov_D\) 表示数据依赖，\(prov_C\) 表示控制依赖。\sys 只在 protected decisions 上约束 \(prov_C\)，例如 tool selection、sink selection、approval scope、policy update 和 delegation target。这样可以表达同一段 PDF text 可以作为 issue body 的 data dependency，却不能作为 GitHub repository 参数的 control dependency。

每个 context cell \(c\) 具有标签：

\[
\ell(c)=\langle origin, integrity, confidentiality, purpose, time, endorsements\rangle .
\]

标签进一步映射到 context authority：

\[
auth(c) \subseteq Modes \times Decisions .
\]

例如，PDF text 可以拥有 \((parameterize, spreadsheet\_cells)\) 和 \((summarize, issue\_body)\)，但不拥有 \((authorize, capability\_request)\) 或 \((sink\_select, github\_repo)\)。实现可以将 \(Modes\) 做成偏序或命名集合；论文中的性质只要求 checker 能判断某个 context 是否被授权影响某个 decision class。

Intent-carrying lease 的抽象形式为：

\[
\kappa =
\langle holder, op, obj, args, I, influence, flow, temporal, budget, expiry, delegation\rangle .
\]

其中 \(I\) 是当前 intent certificate，\(influence\) 约束 control/data provenance，\(flow\) 约束 confidentiality 到 sink 的流向，\(temporal\) 和 \(budget\) 约束事件顺序与调用次数，\(delegation\) 约束可传递给 subagent 的能力集合。

Checker judgment 写作：

\[
\Gamma; I; K; \sigma \vdash e \Rightarrow \sigma' ,
\]

其中 \(\Gamma\) 是全局策略，\(I\) 是 intent certificate，\(K\) 是 active leases，\(\sigma\) 是 runtime state。事件被接受当且仅当存在 \(\kappa \in K\) 使得以下条件同时成立：operation 匹配；object 和 arguments 满足 lease predicate；event purpose 可由 \(I\) 推导；\(prov_C(e)\) 中的每个 context 都拥有影响该 decision class 所需的 authority；\(prov_D(e)\) 满足 flow policy；temporal guard 和 budget 尚未过期；delegated leases 不宽于 parent leases。

由此得到三个安全性质：

\paragraph{Property 1: No Unauthorized Influence.}
对任意 accepted trace \(\tau\)，如果事件 \(e \in \tau\) 属于 protected decision class \(d\)，则对所有 \(c \in prov_C(e)\)，都有 \((mode_d,d) \in auth(c)\)。换言之，未被授权影响 \(d\) 的 context 不能成为该决策的控制依赖。

\paragraph{Property 2: Intent-Bounded Authority.}
对任意 accepted high-impact event \(e\)，存在 active lease \(\kappa\)，且 \(\kappa\) 的 intent predicate 可由 trusted user input、explicit object/sink selection、administrator policy 或 fresh approval 推导。系统不会仅凭 Skill identity、MCP server identity 或工具描述铸造 high-impact authority。

\paragraph{Property 3: Delegation Monotonicity.}
对任意 accepted delegation event \(spawn(a,K_c)\)，child capability set \(K_c\) 是 parent active capability set 的子集或更窄 attenuation。除非 trusted issuer 基于新的 intent certificate 创建 fresh lease，subagent、Skill、MCP tool 或 script 都不能扩大自身权限。

这些性质是 trace-level safety properties。它们依赖于两个工程前提：runtime 必须把 protected decisions 接入 checker，并且 provenance tracker 必须保守地标记事件控制依赖。若某个 OS-level side effect 没有经过任何 adapter，\sys 无法单独观察它；这正是 sandbox 或 ActPlane-style backend 仍然有价值的原因。

\section{实现}

Prototype 由 compiler frontend、deterministic checker 和 runtime adapters 组成。Compiler frontend 使用规则和 LLM-assisted summaries 从 Skills、MCP schemas、local scripts 和 task plan 中推导 candidate effects。所有候选都必须经过 checker。

Runtime adapters 负责把 lease language lowering 到具体边界。Context constructor 将 trusted instructions、untrusted data 和 policy feedback 分区，并记录 provenance。Tool gateway 只暴露 active leases 覆盖的 tools，并验证参数。MCP broker 以 server/tool/method 粒度限制调用，而不是只允许或禁止整个 server。Execution sandbox 将 file/process/network/secrets constraints 下发到底层执行环境。Delegation monitor 确保 subagents 只获得 attenuated leases。

Audit log 记录每次 allow/deny 的 lease、intent、provenance 和 checker reason。这个日志既用于安全解释，也用于 compiler refinement：当 benign action 被拒绝时，runtime 返回结构化 denial，LLM 可以提出更窄的替代计划或请求 fresh approval；checker 仍然决定是否接受。

\section{评估设计}

本文评估围绕四个可证伪问题组织。Run IDs、pilot probes 和 residual analyses 只作为这些问题的证据来源，不作为论文主实验编号。评估的共同原则是固定底层模型、任务顺序、工具描述和 prompt budget，尽量只改变授权层；当某个 baseline 只有 trace-level 或概念性复现时，结果必须标记为 partial comparison，不能写成完整数值胜利。所有 high-impact denials 使用 protected-decision oracle 判定：若某个 protected decision 的 control provenance 来自缺少相应 influence mode 的 context source，则该事件应被拒绝。

\subsection{E1: End-to-End Security and Utility}

E1 检验 \sys 是否在真实 agent workflows 中降低攻击成功率，同时不把 benign task progress 降到不可用。Primary comparison 是同一模型和同一任务集下的 vanilla agent、static ACL、exact-tool ACL、MCP/server allowlist、tool-call guard、OS/sandbox-only wrapper 和完整 \sys。Workloads 覆盖 AgentDojo/InjecAgent-style document or tool-result injection、MCPTox-style MCP poisoning，以及 tau2/tau3 或 MCP utility slices。

主要指标是 attack success、wrong-sink rate、data exfiltration、approval/delegation violations、benign completion、false denial、recovery rate 和 tool errors。Oracle 来自 benchmark security/utility checks 与 \sys protected-decision labels。成功条件是：\sys 相比 access-control baselines 明显降低危险动作，并且 utility loss 有界、可由 denial/recovery 日志解释；若 utility collapse，论文只能主张 \sys 是安全 enforcement layer；若攻击仍通过，则必须定位到缺失的 provenance 或 enforcement boundary。

\subsection{E2: Lease Quality and Authority Reduction}

E2 检验 \sys 生成的 run-time leases 是否比 static extension/server policies 更接近 least privilege。实验单位不是单个工具调用，而是一个 run 在完成任务时暴露的 authority surface。样本来自 security traces、MCP poisoning、stateful utility tasks 和 Skill-style workflows。Oracle 是 blinded expert leases：两名专家在不看 \sys 输出的情况下独立写 minimal leases，经 adjudication 后形成 reference policy。

比较对象包括 domain/global ACL、exact-tool ACL、MCP server allowlist、Skill manifest-style policy、Progent/PCAS/AgentSpec-style policy 和 \sys。指标包括 risk-weighted authority score、exposed tools/methods/sinks/paths/arguments、context-influence breadth、delegation breadth 和 oracle distance。成功条件是 \sys 在保持授权 data use 的同时降低 authority breadth 和 oracle distance；若 symbolic policy baselines 同样窄，则论文应把 least-privilege claim 收窄为 proof-carrying checker operationalization。

\subsection{E3: Context-Influence Mechanism Ablation}

E3 隔离 \sys 的核心机制是否真的是 context influence capability，而不是普通 tool permission 或 taint/provenance tracking。实验比较完整 \sys、no intent certificate、no influence lattice、no control/data split、object-only checker、no-provenance checker，以及 AuthGraph/PACT/AIRGuard-style provenance-authority 和 IFC/taint-style labelers。Workloads 包括 document injection、MCP poisoning、malicious Skill/tool-result 和 residual workflow suites。

关键指标是 protected decisions 中 unauthorized context control 的接受率，以及每个 denial 是否只能由 lease semantics、temporal guards、budget、delegation 或 proof constraints 解释。成功条件是 ablations 或 access-control baselines 接受完整 \sys 拒绝的 protected-decision violations，并且 residual cases 不能被最强 provenance/taint baseline 完全解释。若最强 baseline 已解释所有 denials，论文必须放弃“新的 security boundary”表述，转而主张 \sys 是一种 proof-carrying lease compiler/checker organization。

\subsection{E4: Compiler, Checker, and Recovery Practicality}

E4 检验 LLM-assisted compiler 在 deterministic checker 下是否可用。这里的对象不是模型能力本身，而是一个安全 pipeline：LLM 生成 candidate plans/leases，checker 拒绝 broad/runtime policies 或 missing-proof writes，runtime 返回 structured denial，planner 再生成更窄的替代动作。比较对象包括 LLM-only accept、checker strict、checker+feedback、proof-only/rank-only/filter-only variants 和 checker+CEGAR planner。

指标包括 invalid leases accepted、valid leases rejected、broad/runtime proposals rejected、proof completeness、exact-next precision/recall、denial recovery、dangerous execute、tool errors 和 benchmark task reward。成功条件不是简单增加 executed calls，而是在不激活 broad/runtime authority、零 dangerous execute、零 tool error 的前提下，提高 task-correct executions 或 reward。若 recovery 仍停滞，论文应把 utility claim 收窄为 bounded safety/proof completeness，并将 planner/candidate generation 明确列为 artifact gap。

\section{当前证据与提交前缺口}

当前 prototype 已经支持 checker/gateway、live local callable boundary、benchmark-derived replay、local Qwen/llama.cpp compiler probes、runtime binding、value proof、read/write activation 和 expert-oracle labeling setup。现有证据可以支撑三个谨慎结论：\sys 能在已建模 trace 上阻断 unauthorized context-to-decision influence；run-time leases 明显缩小 object/tool authority surface；deterministic checker 能阻止 LLM-proposed broad/runtime authority 直接执行。

这些结论仍不足以支撑完整顶会 claim。提交前最重要的缺口是：

\begin{itemize}[leftmargin=*]
  \item matched E1 wrapper：同一 workload、同一 model、同一 task set 下与 strong baselines 对比 security 和 utility。
  \item E2 expert oracle：完成 R200 模板的独立专家标注、adjudication 和 oracle-distance scoring。
  \item E3 residual lift：将 lease-semantics residual 从构造 workflow 提升到 benchmark-derived 或 model-loop setting。
  \item E4 planner/CEGAR：证明 recovery 能提高 task-correct executions 或 reward，同时保持零 dangerous execute、零 tool error 和无 broad/runtime lease activation。
\end{itemize}

因此本文当前应写成一个系统安全机制论文，而不是声称已经完成跨 benchmark end-to-end utility。最强的可信 claim 是：\sys 将 context influence 显式编译为 proof-carrying, attenuable capability leases，并在 checker-enforced boundaries 内阻断未授权 context 对 protected decisions 的控制；完整 utility 与 expert-oracle least-privilege 结论需要上述四个主实验闭环。

\section{相关工作定位}

\sys 应被定位为 agent decision authorization model，而不是 extension interface model 或 OS monitor。EIM/bpftime 将扩展能力表示为 host resources，并关注应用扩展的隔离和效率~\cite{eim-bpftime}。\sys 保护的是未来 agent decisions，例如 tool choice、sink selection、authority request 和 delegation。ActPlane 关注如何在 tool layer 之下执行 agent policies~\cite{actplane}；\sys 关注的是这些 policies 应如何从 intent 和 provenance 中生成。SkillGuard/SkillScope 等工作强调 Skill package 的 permission 和 behavior specification~\cite{skillguard}；\sys 是 run-centric 和 cross-extension 的，问的是当前 run 可以允许这个 Skill、这个 MCP server、这份文档和这个 subagent 影响什么。

\bibliographystyle{plain}
\bibliography{intentcap-paper-zh}

\end{document}
